La synthèse qui suit ne résume pas les parties : elle relève ce qui les relie,
ce qu'elles contredisent et ce que l'ensemble laisse ouvert. Elle s'appuie sur
les cinq volets et sur la vérification directe des points contestés.

\subsection{Trois instruments, trois niveaux de garantie}

Le projet assemble trois instruments qui n'offrent pas les mêmes garanties, et
la principale source de confusion est de les traiter comme interchangeables.

\begin{table}[H]
\centering
\caption{Ce que chaque instrument garantit et ne garantit pas.}
\label{tab:instruments}
\small
\begin{tabular}{@{}p{3.4cm}p{4.6cm}p{5.0cm}@{}}
\toprule
Instrument & Garantit & Ne garantit pas \\
\midrule
Frontière d'exécution (\code{executor.py}) & Chaque appel d'outil est journalisé, borné et conforme au périmètre & Une isolation système ; les gardes sont applicatives \\
Bac à sable (\code{sandbox.py}) & Aucun événement Python listé ne survient ; les bras déterministes tournent hors laboratoire & Une isolation native ; aucune mesure de modèle \\
Harnais de calibration & Contrat, oracle séparé, évaluateur qui relit le journal & L'admission de campagne ; la parité exacte des bras \\
\bottomrule
\end{tabular}
\end{table}

Le premier instrument rend l'audit traçable. Le deuxième rend le harnais
portable. Le troisième rend la comparaison possible. Aucun ne prouve qu'un
modèle fait mieux qu'un script, ce qui était l'objet du programme.

\subsection{Les contradictions relevées entre volets}

La confrontation des cinq rédactions a fait apparaître quatre écarts, tous
vérifiés par lecture du code avant d'être retenus ici.

\paragraph{Le seuil du contrôle V0 n'est pas appliqué par la porte.} Le cadre
protocole affirme que la porte applique \code{X = 0,9} ; la partie benchmark
vérifie que la porte se calcule sur « tous succès » et que le contrôle de seuil
est indépendant. Les deux formulations coïncident à trois répétitions mais
divergent au-delà. Le rapport retient la lecture du code.

\paragraph{Le corpus H2 ne code pas le mode d'échec décrit.} Le protocole écrit
que, pour un dossier donné, l'abstention est la seule réponse correcte ; la
lecture de l'oracle montre que celui-ci attend une exposition sur le premier
message et une non-exposition sur le troisième. Le bras par convention converge
d'ailleurs exactement sur l'oracle. Le protocole décrit donc une intention que
le corpus n'encode pas.

\paragraph{Le blocage de l'inférence n'est pas celui qui est communément
allégué.} La condition d'admission est citée comme cause ; la lecture montre
qu'elle est satisfiable en principe, et que la contrainte réelle est
l'absence de familles indépendantes. Ce déplacement importe : il change la
priorité de travail.

\paragraph{La panne de routage du 29 octobre n'est pas consignée.} Elle explique
six heures d'interruption et la reconstruction du conteneur, mais aucun document
du dépôt ne la décrit. C'est une lacune de mémoire du projet, pas un désaccord
entre volets.

\subsection{Un motif commun : le gain absorbé}

Trois épisodes suivent la même forme. Dans la famille H2, le gain du modèle sur
le parseur déclaré est repris par un bras script qui lit par convention. Dans le
bac à sable, le même contraste se reproduit hors ligne, sans modèle. Dans le
rejet de C2, la faille visée se révèle exprimable depuis les entrées publiques,
donc refermable par une automatisation compétente.

Le motif s'énonce ainsi : \emph{ce qui distingue le modèle du script, dans les
familles construites, est une limite de la déclaration plutôt qu'une capacité du
modèle}. Ce n'est pas une démonstration d'infériorité ; c'est l'absence de
situation où la différence apparaît. Or une telle situation est précisément
l'objet du programme.

\subsection{Ce que le projet peut affirmer aujourd'hui}

Quatre affirmations sont défendables, chacune avec sa limite.

Un audit produit des traces vérifiables : chaque appel d'outil est archivé,
référencé et borné, et une preuve perdue interrompt le run plutôt que de
laisser un trou silencieux. La limite est qu'il s'agit de gardes applicatives.

Un harnais peut être exécuté hors du laboratoire, sans broker ni modèle, sous un
garde-fou dont les angles morts sont déclarés. La limite est qu'il ne mesure
rien : il instrumente.

Un troisième bras script absorbe un gain attribué au modèle sur la famille la
plus documentée. La limite est que cela vaut pour huit fixtures et une propriété.

Le corpus est majoritairement détectable par signature, ce qui rend un écart nul
attendu par construction. La limite est que cela concerne le corpus actuel, non
les situations visées.

\subsection{Ce que le projet ne peut pas affirmer}

Aucune supériorité du modèle n'est établie. Aucun intervalle de confiance n'est
calculable. Aucune généralisation n'est possible depuis les mesures existantes.
Et le module d'inférence produit, en l'état, l'apparence d'une mesure qu'il ne
fournit pas : c'est le point le plus sérieux, parce qu'il pourrait être cité de
bonne foi.

\subsection{Les trois décisions qui commandent la suite}

Une première décision porte sur le module d'inférence : le corriger pour qu'il
refuse de produire un chiffre sans familles indépendantes, ou le retirer du
chemin de décision. Tant qu'il existe, chaque nouvelle famille ajoutée augmente
le risque d'une mesure citée hors de sa portée.

Une deuxième décision porte sur la condition d'ouverture : écrire une
automatisation de référence compétente, ou renoncer explicitement aux hypothèses
qui en dépendent. Cette décision rend lisibles plusieurs familles aujourd'hui
indécidables.

Une troisième décision porte sur les familles indépendantes : en produire au
moins trois, ou publier un résultat méthodologique borné en le déclarant tel.
C'est la décision qui détermine si l'inférence pourra un jour exister.

Aucune de ces décisions n'est technique au sens étroit. Toutes portent sur ce
qu'on accepte de ne pas pouvoir dire.

\subsection{Une conclusion lisible}

Le système est construit, instrumenté et documenté. Sa chaîne de preuve est
sérieuse : oracle séparé, évaluateur qui recompte, journal qui ne pardonne pas
une preuve perdue, garde-fou qui déclare ses angles morts. Les défauts relevés
dans ce rapport ont été trouvés en lisant ce que le code fait, non ce qu'il
annonce, et c'est la meilleure preuve que la discipline de vérification
fonctionne.

La question scientifique, elle, reste ouverte. Le programme demandait quand un
modèle fait mieux qu'un script à outils identiques ; les données disponibles ne
répondent pas, et le dispositif actuel ne permettrait pas de répondre même avec
plus de temps machine. Ce constat n'est pas un échec du travail accompli : c'est
le résultat d'avoir regardé les choses en face.

La suite la plus utile n'est pas d'ajouter une famille, mais de rendre
infalsifiable ce qui prétend l'être aujourd'hui.
